% !TeX program = xelatex
% Overleaf: Menu -> Compiler -> XeLaTeX
% 本文件严格映射 REPORT_OUTLINE.md：只保留报告大纲与写作提示。
\documentclass[UTF8,11pt,a4paper,fontset=none]{ctexart}

\usepackage{fontspec}
\usepackage{xeCJK}
\IfFontExistsTF{Noto Serif CJK SC}{
  \setCJKmainfont{Noto Serif CJK SC}
  \setCJKsansfont{Noto Sans CJK SC}
  \setCJKmonofont{Noto Sans CJK SC}
}{
  \setCJKmainfont{Songti SC}
  \setCJKsansfont{PingFang SC}
  \setCJKmonofont{PingFang SC}
}

\usepackage[a4paper,margin=2.5cm]{geometry}
\usepackage{amsmath,amssymb,bm}
\usepackage{booktabs,tabularx,array,multirow,longtable}
\usepackage{graphicx,xcolor}
\usepackage{enumitem}
\usepackage{caption,subcaption,float}
\usepackage{microtype}
\usepackage{hyperref}
\usepackage[nameinlink,noabbrev]{cleveref}
\usepackage[numbers,sort&compress]{natbib}

\definecolor{linkblue}{HTML}{175A9C}
\definecolor{promptblue}{HTML}{245A8D}
\hypersetup{
  colorlinks=true,
  linkcolor=linkblue,
  citecolor=linkblue,
  urlcolor=linkblue,
  pdftitle={Bomberman Agent 项目报告},
  pdfauthor={作者姓名}
}

\setlength{\parindent}{2em}
\setlength{\parskip}{0.25em}
\setlist{nosep,leftmargin=2em}
\captionsetup{font=small,labelfont=bf}

% 课程要求：每个章节或小节标题后标明主要作者。
\newcommand{\primaryauthor}[1]{\texorpdfstring{\enspace\normalfont\small[主要作者：#1]}{ [主要作者：#1]}}
\newcommand{\authorsection}[2]{\section{#1\primaryauthor{#2}}}
\newcommand{\authorsubsection}[2]{\subsection{#1\primaryauthor{#2}}}
\newcommand{\authorsubsubsection}[2]{\subsubsection{#1\primaryauthor{#2}}}
\newcommand{\pendingauthor}{待填写}

% 写作提示开关：提交前改为 \draftguidancefalse，所有蓝色提示会消失。
\newif\ifdraftguidance
\draftguidancetrue
\newcommand{\writingnote}[1]{%
  \ifdraftguidance
    \par\noindent\begingroup\color{promptblue}\small
    \fbox{\begin{minipage}{0.94\linewidth}\textbf{写作提示：}#1\end{minipage}}%
    \par\endgroup
  \fi
}
\newcommand{\todoitem}[1]{%
  \ifdraftguidance\textcolor{promptblue}{\textbf{[待补：#1]}}\fi
}

\newcommand{\experimentcardtitle}[1]{\par\medskip\noindent\textbf{#1}\par\nobreak\smallskip} \newcommand{\experimentcardfield}[2]{\noindent\textbf{#1：}#2\par} \title{\todoitem{报告标题}}
\author{
  \todoitem{成员姓名 1}\quad
  \todoitem{成员姓名 2}\quad
  \todoitem{成员姓名 3}\\
  \small Machine Learning Essentials, Summer Semester 2026
}
\date{\today}

\begin{document}
\maketitle

\writingnote{提交前确认课程是否接受中文及其字数折算规则。课程要求的总篇幅约为每名成员 4,000 English words，且不应明显超出。报告中不要使用大学 Logo。}

\begin{center}
  \textbf{公开代码仓库：}\todoitem{公开仓库 URL}\\
  \textbf{额外运行依赖：}\todoitem{没有则写“无”；有则列名称与版本}
\end{center}

\begin{abstract}
\writingnote{用一段自洽文字概括大纲中的研究问题、三条模型路线、统一冻结评估方法、最重要的定量结果，以及为什么最终选择 B33。最后写清结论只适用于已运行协议；不要引用文献，也不要加入正文未报告的结果。}
\todoitem{摘要正文}
\end{abstract}

\noindent\textbf{关键词：}\todoitem{3--5 个关键词}

\tableofcontents
\clearpage

\writingnote{提交检查：摘要和 Background 使用完整正文，不保留提示或示意图占位；报告开头列出最终 Agent 的 NumPy 与 PyTorch 版本；所有标题填写真实主要作者；所有数字可追溯到 report-assets/tables/*.csv 或其 source\_path；not\_reported 和 unrun 不转换为 0；报告 PDF 不上传公开仓库；团队成员须理解、核验并以自己的表达修订 AI 辅助初稿。}

\authorsection{引言}{\pendingauthor}

\authorsubsection{问题与研究动机}{\pendingauthor}
\writingnote{说明稀疏反馈、延迟炸弹信用、目标重选和对手不确定性；强调训练 reward 不等于冻结游戏表现，得分、安全和 0.5 秒 CPU 时限共同决定选模。}
\todoitem{正文}

\authorsubsection{研究问题}{\pendingauthor}
\writingnote{依次回答：不同价值学习方法能否学会金币导航并迁移到炸箱和对战；等待、循环、自杀和目标丢失分别能由哪些特征、奖励和生存约束改善；三条路线的主要瓶颈及哪些改进有效、无效或证据不足；在得分、生存、推理时间和可复现性之间为什么最终选择 B33。}
\todoitem{正文}

\authorsection{背景}{\pendingauthor}
\writingnote{描述 Bomberman 环境与 Task 1--4；概述 Q-learning、Double Q、DQN、Double DQN、dueling、PER、n-step 与蒸馏；介绍离散、tile-coded 连续、MLP 连续和 17 通道空间表示；说明奖励塑形、课程学习与安全动作准入。}
\todoitem{正文}

\authorsection{项目规划}{\pendingauthor}

\authorsubsection{团队结构与协作}{\pendingauthor}
\writingnote{路线一为 Continuous Double DQN Task 1--4 与 survival mask；路线二为 Rainbow Lite Task 1--4、feature/reward 与延迟信用；路线三为 Q-learning/CNN Task 1--2 与跨学习器迁移。说明三条路线共享代码、指标、失败案例和选模规则，不得描述为三个独立项目。}
\todoitem{正文}

\authorsubsection{门槛与证据管理}{\pendingauthor}
\writingnote{说明“训练→开发选模→独立确认→主验证→打包”的流程；区分 training seed 与 environment seed，未运行记为 unrun，不可记为零；冻结 checkpoint、关闭探索、使用单线程 CPU，并记录配置/checkpoint 哈希和 paired bootstrap。}
\todoitem{正文}

\authorsection{方法}{\pendingauthor}

\authorsubsection{共享状态与特征}{\pendingauthor}
\writingnote{说明离散键、84 维 tile coding、70/78/84 维连续表示、117 维 phase 表示和 17 通道 CNN；将 Feature v1--v11 按导航、历史/循环、危险、资源、对手和目标可达性归类；明确版本号仅作复现标识，不代表单调性能顺序。}
\todoitem{正文}

\authorsubsection{共享奖励与生存准入}{\pendingauthor}
\writingnote{将 Reward r1--r22 按正式事件、势函数、WAIT/循环、炸弹信用、生存和对手事件归类；Feature、reward、mask 是独立版本轴；区分 physical legal action、finite-horizon survivable action、veto-only mask 和 fallback；说明 Survival-mask v1--v9 由具体反例驱动，并保留 world 27155 限制。}
\todoitem{正文}

\authorsubsection{模型与统一评估}{\pendingauthor}
\writingnote{介绍表格 Q/Double Q/Watkins Double Q($\lambda$)、Continuous/CNN Double DQN 与蒸馏 CNN；Rainbow Lite 定义为 dueling + proportional PER + fixed 4-step，不称完整 Rainbow；指标和可比性严格遵循 report-assets/plan/experiment-protocol.md。}
\todoitem{正文}

\writingnote{图表占位：Figure 1 为课程训练与证据流程；Figure 2 为 Die Hardest 推理和动作准入；Table 1 为模型族及证据状态。插入正式图表时使用 report-assets 中的可追溯资产。}
\todoitem{Figure 1、Figure 2 与 Table 1}

\authorsection{训练}{\pendingauthor}
\writingnote{说明 Task 1 导航、Task 2 炸箱与隐藏金币、Task 3 弱对手、Task 4 三个规则对手；当前任务、旧任务保留、安全和运行时共同决定晋级；严格 resume 与 policy-only transfer 分开，合同变化时显式迁移或重训；表格/MLP 主要使用 CPU，CNN 训练使用 GPU，最终评估为 CPU-only；保留 CNN 输入 bug、WAIT/循环、Task 2 自杀、Task 3 计数缺陷及 Task 4 超时/中断等失败记录。}
\todoitem{正文}

\authorsection{实验与结果}{\pendingauthor}
\label{sec:results}
\writingnote{本章是报告重点。每个实验段落必须回答：问题、证据、诊断、解决方案、验证、结果、有限答案和下一决策。}

% REPORT_OUTLINE.md 明确将总体逻辑编号为 6.0。
\setcounter{subsection}{-1}
\authorsubsection{总体实验逻辑与比较边界}{\pendingauthor}
\writingnote{Task 1--4 依次隔离导航、炸箱逃生、弱对手竞争和强对手压力，后续任务必须回测早期能力；三条路线共享 Feature、Reward、Survival Mask 与冻结评估，但针对表示、信用和安全提出不同假设；课程 agents 作为兼容性、课程或压力测试对象，缺少同协议冻结结果时不建立 baseline 数值排名；verified\_raw 支持正式差值，verified\_summary 描述历史结果，not\_comparable 仅作诊断，unrun 明确为未运行。}
\todoitem{正文}

\authorsubsection{Continuous Double DQN：从导航到可控生存}{\pendingauthor}

\authorsubsubsection{Task 1 $\rightarrow$ Task 2：为什么会炸箱仍会自杀？}{\pendingauthor}
\writingnote{按“问题—证据—诊断—解决—验证—结论—决策”写作。早期候选为 5.45 coins、76.65 crates、30\% suicide；解决方案是将 survival mask 用于行为、推理和 bootstrap；主验证为 7.25 coins、100.26 crates、0\% suicide、100\% Task 1 retention；结论限定为安全准入解决晋级阻塞，但不等于学会全部长期规划。}
\todoitem{正文}

\authorsubsubsection{Task 3：能否改善竞争并保留旧能力？}{\pendingauthor}
\writingnote{按统一七步叙事写作。报告三 seed confirmation 与 seed22/c150 main validation；三条独立训练链得分增量为 +1.92、+1.23、+2.56，均值 +1.90，样本标准差 0.67；每条链的 100 个世界是环境重复，不写成 300 次独立训练；得分 5.91 $\rightarrow$ 7.14，差值 1.23，CI [0.27, 2.14]；第一名率 41\% $\rightarrow$ 53\%；击杀 0.45 $\rightarrow$ 0.44，不称击杀改善。}
\todoitem{正文与 Table 2：Task 3 父子配对}

\authorsubsubsection{Task 4：专项训练为何没有稳定收益？}{\pendingauthor}
\writingnote{按统一七步叙事写作。说明 Shared-parent admission failure；将 E1/E2/E3、Frozen C/S 和 score campaign 分开呈现；结论只限定于已运行协议。}
\todoitem{正文}

\authorsubsection{Rainbow Lite：从奖励塑形到安全动作约束}{祝雨嫣}

Bomberman 的困难不只在于选择眼前动作，还在于把数步以后才出现的炸箱、金币、击杀或死亡结果归因给先前决策。Rainbow Lite 在 Double DQN 上加入 dueling network、比例优先经验回放和四步回报，分别改善状态价值与动作优势的分解、关键样本复用以及延迟信用传播。它没有采用完整 Rainbow 的分布式价值学习和 noisy network，因此保留 ``Lite'' 的称呼。

理解这条路线需要先区分三个容易混在一起的部件。feature 决定模型能看到什么，例如资源方向、近期动作和对手位置；reward 决定训练时鼓励什么，例如收集资源、减少停滞或主动攻击；action mask 则在模型给出 Q 值以后，直接排除当前不允许执行的动作。前两者改变学习问题，最后一项改变实际可选动作集合。因而，“换 reward 后自杀减少”和“mask 阻止了自杀动作”是性质不同的证据，不能写成同一种改进。

早期开发主要在 mask-v1 下反复增加 feature 和 reward。它们有时提高某个小开发集的得分，也有时减少 WAIT，但自杀并未稳定消失；当表示和奖励同时变化时，也无法判断究竟是哪一项起作用。这正是后续实验不再追逐单个 checkpoint 高点，而改为依次回答三个更简单的问题的原因：先验证 r18 能否针对 WAIT，再检验更复杂的 V6+r19 是否值得，最后把训练 mask 与推理 mask 拆开。这样的顺序使每一步都只承担有限的结论。

课程本身也逐步增加难度。Task 1 只要求在固定场景中收集金币；Task 2 加入炸箱和资源搜索；Task 3 增加一名对手；Task 4 则面对三名对手。前一任务的 checkpoint 会成为后一任务的起点，所以后期得分不只由当前配置决定，也携带父模型的行为习惯。一个模型可能在 Task 4 更积极，却忘记 Task 1 的导航；也可能通过安全约束避免死亡，却因为频繁 WAIT 而失去竞争力。因此，最终选择必须同时检查当前任务、早期任务保留和行为失败模式。

本路线最后得到一个比“哪个 checkpoint 分数最高”更重要的结论：在固定训练 seed 和预算下，r18 比 r17 更适合处理长期 WAIT；继续增加 feature/reward 复杂度没有形成稳定收益；mask-v5 则明显改善了危险动作控制。基于下面四项核心实验，最终保留 V5+r18+mask-v5。是否还应从 Task 1 重训完整课程链是另一项决策，不能由 mask 的安全收益直接推出，因此在后文单独比较旧 anchor 与完整重训路线。这里的“最终”只表示该组件组合是受控候选中性能、安全与复杂度较平衡的方案，并不表示已经证明其为所有配置中的全局最优。

为了避免把开发过程写成版本流水账，本节只围绕三个逐步收窄的问题展开：停滞应由 reward 还是更复杂的表示来解决，安全性应由策略自己学习还是由 action mask 约束，以及进入 Task 4 时是否值得从 Task 1 重训完整课程链。所有新增训练最多运行 500 局，只比较 c500；冻结评估关闭探索，在相同的 100 个 environment worlds 上逐 world 配对，并用 10,000 次 paired bootstrap 给出 95\% 置信区间。environment worlds 只刻画地图与对手随机性，不能替代独立训练 seed。

这里的配对含义是：两个候选面对完全相同的地图、初始随机性和对手，再计算同一 world 内的指标差。这样可以减少“一个模型碰巧遇到更容易地图”带来的噪声。若 95\% 置信区间跨过 0，本文只说观察到方向或数值差，不把它写成明确改善；若比较只有一个训练 seed，也只讨论该 seed 下的结果。训练 reward 和 loss 不作为最终选模指标，因为它们可能上升而实际游戏得分、安全性或行为质量没有改善。

\authorsubsubsection{r18 是否比 r17 更适合处理长期 WAIT？}{祝雨嫣}

结果先说：在这组受控比较中，r18 确实消除了 r17 仍然出现的长期 WAIT，而且没有观察到得分退化。两条分支从同一个 Task 1 父权重出发，都使用 V5、mask-v1、训练 seed 11 和 500 局预算；唯一变化是 reward 从 r17 换为 r18。随后在相同的 100 个 worlds 上进行冻结配对评估。

\begin{table}[!htbp]
\centering
\caption{Task 2 中 r18 相对 r17 的冻结配对结果。差值为 r18$-$r17。}
\label{tab:rainbow-lite-r18-r17}
\small
\begin{tabularx}{\linewidth}{l *{3}{>{\centering\arraybackslash}X} >{\centering\arraybackslash}p{2.5cm}}
\toprule
指标 & r17 & r18 & 差值 & 95\% CI \\
\midrule
长 WAIT & 10\% & 0\% & $-10$ pp & [$-16$,$-5$] \\
得分/金币 & 8.21 & 8.40 & $+0.19$ & [$-0.04$,$0.41$] \\
炸箱数 & 114.48 & 114.72 & $+0.24$ & [$-1.26$,$1.66$] \\
放置炸弹 & 40.03 & 38.08 & $-1.95$ & [$-2.75$,$-1.17$] \\
自杀率 & 0\% & 0\% & 0 pp & [0,0] \\
\bottomrule
\end{tabularx}
\end{table}

r18 将长期 WAIT 从 10\% 降至 0\%，区间不跨 0；得分和炸箱数则基本保持。炸弹数量显著减少但炸箱数不变，说明它更可能减少了无效动作，而不是简单压低行动频率。因此，“r18 在不损害主要任务指标的情况下比 r17 更适合处理 WAIT”这一假设在固定 seed 11 下得到\textbf{支持}。这仍是单训练 seed 结论，不能外推为跨 seed 的稳定优势。


\authorsubsubsection{继续增加 feature/reward 复杂度是否有益？}{祝雨嫣}


团队首先把当时更复杂的 V6+r19 作为一个整体 bundle，与较简单的 V5+r18 进行正式比较。为避免把不同协议的局部高点混作公平排名，实验只比较两条共享 Task 3 父权重、mask-v1、训练 seed 11 和 500 局预算的分支，并在相同的 100 个 Task 4 worlds 上评估。




结果没有支持复杂方案成为更好的最终选择。V6+r19 的确把长期 WAIT 从 35\% 降到 1\%，但自杀率同时由 35\% 升到 57\%；平均得分仅增加 0.20，置信区间跨过 0。

\begin{table}[!htbp]
\centering
\caption{Task 4 中复杂 bundle 的冻结配对结果。差值为 V6+r19$-$V5+r18。}
\label{tab:rainbow-lite-v6-v5}
\small
\begin{tabularx}{\linewidth}{l *{3}{>{\centering\arraybackslash}X} >{\centering\arraybackslash}p{2.5cm}}
\toprule
指标 & V5+r18 & V6+r19 & 差值 & 95\% CI \\
\midrule
得分 & 2.47 & 2.67 & $+0.20$ & [$-0.51$,$0.91$] \\
自杀率 & 35\% & 57\% & $+22$ pp & [$+9$,$+35$] \\
存活率 & 59\% & 37\% & $-22$ pp & [$-35$,$-9$] \\
第一名率 & 20\% & 22\% & $+2$ pp & [$-10$,$+14$] \\
长 WAIT & 35\% & 1\% & $-34$ pp & [$-44$,$-25$] \\
\bottomrule
\end{tabularx}
\end{table}

这不是“复杂特征完全无用”，而是复杂 bundle 把失败模式从停滞转成了更频繁的致命动作，且没有形成可靠的得分或排名收益。因此，“继续增加 feature/reward 复杂度会带来更好的综合表现”这一假设\textbf{不受支持}。在当前预算下，继续承担 V6+r19 的实现成本缺乏依据。

在这项正式对照未显示稳定综合收益后，团队仍继续检验其他复杂化方向：V6 Stable 调整对手压力与稳定性；V7 增加击杀状态并让阶段奖励可观察；V8 加入可生存的炸箱机会；V9 删除常量和冗余字段、压缩表示；V10 与 V11 分别尝试区域密度和可达箱目标；Spatial V6 则改用空间棋盘表示。在 r19 之后，r20--r22 又进一步调整攻击信用与阶段奖励，包括取消预测性 BOMB 正奖励，以及按资源和首杀状态改变势能。\par 这些后续分支有时在小样本或单个 checkpoint 上出现积极信号，但没有在统一协议下确认稳定优于 V5+r18；相反，它们还伴随更高自杀率、过度攻击、checkpoint 大幅波动或长期 WAIT。由于各分支的父模型、reward、训练预算和评估规模未全部保持一致，这些结果只能作为诊断，不能构成 V6--V11 或 Spatial V6 的公平排名，也没有推翻前述正式配对的判断。\par 综合而言，现有证据不是证明所有额外特征都无效，而是尚未确认继续增加 feature/reward 复杂度能在当前预算下稳定优于 V5+r18；其额外实现和调试成本因此缺乏依据。

\authorsubsubsection{mask-v5 的收益来自推理约束还是重新训练？}{祝雨嫣}

冻结同一权重只切换 mask 的先导实验已经给出一个清楚但不完整的信号：Task 3 自杀率从 mask-v1 的 31\% 降为 mask-v5 的 0\%（差值 $-31$ 个百分点，95\% CI [$-40$,$-22$]），得分差 $+0.28$ 的区间却跨过 0。它说明 mask-v5 能在推理时阻止一批危险动作，但不能回答模型是否还需要在新的动作约束下重新学习。

因此，正式实验从同一个 V5+r18 Task 2 父模型出发，以训练 seeds 11、22、33 分别训练 mask-v1 和 mask-v5 六条 Task 3 分支；每条只训练 500 局。随后每个 checkpoint 又分别在 eval-mask-v1 与 eval-mask-v5 下运行相同的 100 个 worlds，形成训练 mask 与推理 mask 的 $2\times2$ 对照。这个设计把“推理时临时换 mask 已经足够吗”和“在 mask-v5 下重新训练是否有额外价值”拆成两个可检验的问题。

\begin{table}[!htbp]
\centering
\caption{三条下游训练分支的最终契约比较。每个单元依次报告得分/自杀率/长 WAIT。}
\label{tab:rainbow-lite-mask-contract}
\small
\begin{tabularx}{\linewidth}{c >{\centering\arraybackslash}X >{\centering\arraybackslash}X >{\centering\arraybackslash}p{3.1cm}}
\toprule
训练 seed & \shortstack{训练 mask-v1\\评估 mask-v1} & \shortstack{训练 mask-v5\\评估 mask-v5} & 得分差的 95\% CI \\
\midrule
11 & 6.03/20\%/53\% & 6.32/0\%/72\% & $+0.29$ [$-0.70$,$1.34$] \\
22 & 5.58/24\%/48\% & 6.81/0\%/33\% & $+1.23$ [$0.33$,$2.13$] \\
33 & 6.20/21\%/58\% & 7.34/0\%/70\% & $+1.14$ [$0.21$,$2.08$] \\
\bottomrule
\end{tabularx}
\end{table}

最终契约在三个 seed 上都把自杀率降至 0\%，而基线分别为 20\%、24\% 和 21\%；得分差方向也一致为正，其中 seeds 22、33 的区间不跨 0。不过，长 WAIT 的变化为 $+19$、$-15$ 和 $+12$ 个百分点，并不稳定。进一步拆开 $2\times2$ 对照后，只在推理时从 v1 切换到 v5，就能在所有六种“训练权重$\times$seed”组合中把 20\%--26\% 的自杀率降至 0\%；但在 v5 下重新训练相对临时换 mask 的得分差为 $-0.20$、$+0.88$ 和 $+1.00$，并非三个 seed 都受益。

因此，该实验\textbf{部分支持}最终选择：mask-v5 的安全收益清楚且跨 seed 一致，最终契约的得分也未出现方向性退化；但额外收益主要来自推理约束，重新训练带来的性能增量尚不稳定，而且 WAIT 仍是代价。它不能被写成“mask-v5 同时稳定提高所有指标”。

三个训练 seed 是共享同一 Task 2 父模型后的下游重复，而不是三条从零开始的完整课程链。因此，seed 间方向一致性可以提高对 Task 3 结论的信心，但仍不能代替完整多 seed 课程复现。

\authorsubsubsection{完整重训是否优于沿用旧 anchor？}{祝雨嫣}

这里比较的不是两种网络结构，而是两种进入 Task 4 的训练历史。旧 anchor 路线沿用早期 Task 3 checkpoint c350；完整重训路线则用固定的 V5+r18+mask-v5 从 Task 1 重新走到 Task 3 c300。首先冻结这两个“进入 Task 4 前的模型”，比较它们在 Task 1--3 的保留能力；随后只加载各自的 policy weights，重置 optimizer、replay、epsilon、随机数和阶段状态，再用相同的 seed 11、V5+r18+mask-v5 与 500 局预算训练 Task 4。这样得到的是两条路线各自“完成 Task 4 后的模型”。

\begin{table}[!htbp]
\centering
\caption{旧 anchor 与完整重训路线在进入和完成 Task 4 训练时的冻结结果。旧与重训分别单列；每项均为 100 个配对 worlds。}
\label{tab:rainbow-lite-retrain-chain}
\scriptsize
\begin{tabularx}{\linewidth}{p{1.85cm}l *{6}{>{\centering\arraybackslash}X}}
\toprule
模型阶段 & 任务 & \multicolumn{2}{c}{得分} & \multicolumn{2}{c}{长 WAIT} & \multicolumn{2}{c}{存活率} \\& & 旧 & 重训 & 旧 & 重训 & 旧 & 重训 \\
\midrule
进入 Task 4 前 & Task 1 & 21.47 & 17.55 & 0\% & 4\% & 100\% & 100\% \\
进入 Task 4 前 & Task 2 & 7.25 & 7.76 & 1\% & 0\% & 100\% & 100\% \\
进入 Task 4 前 & Task 3 & 7.04 & 7.75 & 7\% & 0\% & 100\% & 100\% \\
\midrule
完成 Task 4 后 & Task 1 & 43.51 & 31.72 & 0\% & 37\% & 100\% & 100\% \\
完成 Task 4 后 & Task 2 & 4.94 & 5.57 & 35\% & 6\% & 100\% & 100\% \\
完成 Task 4 后 & Task 3 & 5.73 & 5.89 & 51\% & 54\% & 100\% & 100\% \\
完成 Task 4 后 & Task 4 & 2.83 & 3.21 & 38\% & 11\% & 68\% & 85\% \\
\bottomrule
\end{tabularx}
\end{table}

表中的 Task 1 数字不是最初 Task 1 checkpoint 的成绩，而是后续模型对 Task 1 的回测。Task 1 c0300 的 50.00 来自 20 个开发 worlds；本表使用另一组 100 个冻结 worlds，且比较的是已经进入 Task 3 或完成 Task 4 的权重，因此二者不能直接相减并解释为遗忘量。表内可以严格比较的是同一行的两条路线。\par 结果不支持“完整重训全面更强”。进入 Task 4 前，完整重训路线在 Task 2 得分上高 0.51（95\% CI [$0.15$,$0.87$]），Task 3 的 $+0.71$ 差值区间跨 0，但 Task 1 已低 3.92。完成相同 Task 4 训练后，完整重训路线在 Task 2 得分、Task 4 存活率和长 WAIT 上更好；旧 anchor 路线则明显保留了更多 Task 1 能力。完整重训路线的 Task 1 得分低 11.79（95\% CI [$-13.35$,$-10.27$]），并出现 37\% 长 WAIT。因此，“完整重训会同时改善早期任务保留且不损害 Task 4”这一假设\textbf{不受支持}。

路线选择取决于目标。若优先均衡保留 Task 1--4，现有结果更偏向旧 anchor；若最终模型主要面向 Task 4，并把存活和避免长期 WAIT 放在更高优先级，则完整重训路线更合适。本研究采用后一标准，因此保留完整重训链作为 Rainbow Lite 的 Task 4 候选：它在该配对中将存活率从 68\% 提高到 85\%，并将长 WAIT 从 38\% 降到 11\%。这是一项面向 Task 4 的工程取舍，不是完整重训整体性能最优的证明。

这个实验检验的是进入 Task 4 前的训练历史是否会在控制住后续训练条件后继续影响保留与迁移。结果确认这种影响确实存在，但方向随任务变化；同时它仍只是 seed 11 的课程链比较，不能写成多 seed 稳定复现。因此，后文直接汇总最终选择及其证据边界。

\authorsubsubsection{最终选择与证据边界}{祝雨嫣}

四项核心实验把开发过程收敛为三个确定发现。第一，r18 在单 seed 配对实验中消除了长期 WAIT，且没有牺牲 Task 2 的资源指标。第二，更复杂的 V6+r19 虽能减少 WAIT，却显著增加自杀，没有稳定综合收益。第三，mask-v5 在三个下游训练 seed 上都稳定消除了自杀，但其得分收益和 WAIT 代价并不统一。相应的判定依次为：实验 1 \textbf{支持}，实验 2 \textbf{不支持}，实验 3 \textbf{部分支持}，实验 4 \textbf{不支持}“完整重训全面更强”。

这些证据共同限定了最终选择的含义。V5+r18+mask-v5 不是一次把所有问题解决的复杂方案，而是主动停止继续堆叠 feature/reward 后得到的折中：r18 针对 WAIT，V5 保留足够的资源方向信息，mask-v5 负责阻止无法证明安全的动作，完整重训路线则是在优先 Task 4 安全与训练契约一致性时采用的工程选择，而不是全任务最优结论。

从工程决策看，这一结论也意味着后续工作不应继续盲目扩充输入和奖励项。更直接的方向是改进 mask-v5 在“安全但可行动”场景中的放行规则，并加入针对 Task 1 遗忘的保留机制；这两个问题都能由现有冻结协议直接验证，而无需再同时改变多个组件。

最后，新增实验中只有 mask 对照使用三个下游训练 seed；r17/r18、V5/V6 和完整课程链比较仍固定 seed 11。world-level bootstrap 只量化环境随机性，不量化重新训练带来的方差。故本研究可以说明为什么在当前预算和候选集合中选择 V5+r18+mask-v5，却不能声称某一个 feature、reward 或 mask 独立解释了全部性能变化，也不能声称该组合在任意训练 seed 上都占优。

\authorsubsection{Q-learning 与 CNN：共享成果能否跨模型迁移}{\pendingauthor}

\authorsubsubsection{Task 1：两种替代学习器能否导航？}{\pendingauthor}
\writingnote{按统一七步叙事写作。Double Q($\lambda$) 为 48.95 coins、96\% all-coins；Distilled CNN 为 49.84 coins、96\% all-coins；额外 TD fine-tuning 退化到 37.21 coins、55\% all-coins。}
\todoitem{正文}

\authorsubsubsection{Task 2：安全炸箱为何仍不能完成金币目标？}{\pendingauthor}
\writingnote{按统一七步叙事写作。Q r20 为 3.95/9 coins、53.45 crates、0\% suicide、75\% long ping-pong；CNN D02 为 2.60/9 coins、44.15 crates、0\% suicide、0\% all-coins；共同瓶颈是箱后长期目标重选；两者未通过 Task 2，因此停止进入 Task 3/4。}
\todoitem{正文与 Table 3：Q/CNN Task 1--2 代表结果}

\authorsubsection{跨路线综合与最终选择}{\pendingauthor}
\writingnote{按统一七步叙事写作。用表格总结三条路线的原始问题、方案、证据、结论和决策；不建立跨协议统一排行榜；将 B33 的 1,000-world benchmark 与 Die Hardest 的六局 package equivalence 分开；最终选择依据为完整 Task 1--4 谱系、4.231 历史均分、50.2\% 含并列第一、94.7\% survival 和可复现打包；明确 world 27155、Docker/官方硬件未验证以及该选择不是普遍最优证明。}
\todoitem{正文与 Table 4：三条路线问题—解决—结论综合}

\authorsection{结论}{\pendingauthor}
\writingnote{总结：多种价值学习器可完成 Task 1，但迁移深度不同；历史/reward 改变停滞，mask 减少自杀，二者都不自动解决长期规划；Task 3 得分改善，而 Task 4 增量方案缺少稳定收益；B33 是证据完整性与工程约束下的最终选择。不得引入正文未报告的新实验。}
\todoitem{正文}

\bibliographystyle{plainnat}
\bibliography{references}

\appendix
\authorsection{证据索引与复现信息}{\pendingauthor}
\writingnote{仅放置正文容纳不下的证据索引和复现信息；附录不得替代正文中的核心方法、实验结果与最终选择依据。}
\todoitem{补充内容}

\end{document}
